\section*{Hoofdstuk \ref{ch:b1:rate-laws} --- Chemische kinetiek: snelheidswetten}

\begin{solution}{exo:b1:rate-laws:1}
$v = \frac14 \times 2.0 \times 10^{-4} = \qty{5.0e-5}{mol.L^{-1}.s^{-1}}$;
dizuurstof wordt met dezelfde snelheid gevormd, \qty{5.0e-5}{};
\ce{N2O5} verdwijnt met
$2v = \qty{1.0e-4}{mol.L^{-1}.s^{-1}}$.
\end{solution}

\begin{solution}{exo:b1:rate-laws:2}
Orde 0: \unit{mol.L^{-1}.s^{-1}}; orde 1: \unit{s^{-1}}; orde 2:
\unit{L.mol^{-1}.s^{-1}}; orde 3: \unit{L^2.mol^{-2}.s^{-1}}.
\end{solution}

\begin{solution}{exo:b1:rate-laws:3}
$t_{1/2} = \ln2/k = 0.693/(3.0 \times 10^{-3}) = \qty{231}{s}$. Na
\qty{600}{s} is $\eu^{-kt} = \eu^{-1.8} = 0.165$: er blijft 16.5\,\%
over.
\end{solution}

\begin{solution}{exo:b1:rate-laws:4}
Een \omterm{def:b1:rate-laws:half-life}{halveringstijd} die evenredig is met $1/a_0$ betekent orde 2. Als het
halveren van $a_0$ de \omterm{def:b1:rate-laws:half-life}{halveringstijd} halveert, is $t_{1/2} \propto a_0$:
orde 0.
\end{solution}

\begin{solution}{exo:b1:rate-laws:5}
$[\ce{A}]$ is niet lineair in $t$ (verliezen 0.0259, 0.0192, 0.0142,
0.0106). $\ln[\ce{A}] = -2.303$, $-2.602$, $-2.902$, $-3.202$, $-3.503$:
het daalt met 0.300 per \qty{10}{min}, een rechte. Orde 1, $k =
\qty{0.0300}{min^{-1}}$.
\end{solution}

\begin{solution}{exo:b1:rate-laws:6}
Een verdubbeling van $[\ce{A}]_0$ vermenigvuldigt $v_0$ met $4 = 2^2$:
orde 2 in \ce{A}. Een verdubbeling van $[\ce{B}]_0$ vermenigvuldigt ze
met 2: orde 1 in \ce{B}. $v =
k[\ce{A}]^2[\ce{B}]$ met $k = 2.0 \times 10^{-6}/(0.010^2 \times 0.010) =
\qty{2.0}{L^2.mol^{-2}.s^{-1}}$.
\end{solution}

\begin{solution}{exo:b1:rate-laws:7}
\ce{B} is in honderdvoudige overmaat: zijn concentratie blijft
\qty{0.50}{mol/L}, en $v = k'[\ce{A}]$ met $k' = k[\ce{B}]_0$. Dus
$k = 0.012/0.50 = \qty{0.024}{L.mol^{-1}.s^{-1}}$.
\end{solution}

\begin{solution}{exo:b1:rate-laws:8}
$E_a = R\ln5/(1/300 - 1/320) = 8.314 \times 1.609/(2.083 \times 10^{-4}) =
\qty{64.2}{kJ/mol}$. Bij \qty{310}{K}: $\ln k = \ln(1.0 \times 10^{-3}) +
\frac{E_a}{R}\big(\frac{1}{300} - \frac{1}{310}\big) = -6.908 + 0.831$,
$k = \qty{2.3e-3}{s^{-1}}$.
\end{solution}

\begin{solution}{exo:b1:rate-laws:9}
Hier is $a = 2$: $1/[\ce{A}] = 1/a_0 + 2kt$, $t_{1/2} = 1/(2k a_0) = 1/(2
\times 0.50 \times 0.020) = \qty{50}{s}$. Bij \qty{100}{s}: $1/[\ce{A}] = 50
+ 100 = \qty{150}{L/mol}$, $[\ce{A}] = \qty{6.7e-3}{mol/L}$.
\end{solution}

\begin{solution}{exo:b1:rate-laws:10}
$[\ce{A}] = a_0(1 - f) = a_0\eu^{-akt_f}$ geeft $t_f = -\ln(1-f)/(ak)$, en
in \omterm{def:b1:rate-laws:half-life}{halveringstijden} $t_f/t_{1/2} = -\ln(1-f)/\ln2$: 1 voor 50\,\%, 2 voor
75\,\%, 3.32 voor 90\,\%, 9.97 voor 99.9\,\%: ongeveer tien
\omterm{def:b1:rate-laws:half-life}{halveringstijden} voordat de reactie ``af'' is.
\end{solution}

\begin{solution}{exo:b1:rate-laws:11}
$m(t) = 10 - 0.40\,t$ (\unit{mg}, \unit{h}): orde 0. \omterm{def:b1:rate-laws:half-life}{Halveringstijd}
$10/(2 \times 0.40) = \qty{12.5}{h}$; leeg na \qty{25}{h}. De dosis die
per uur wordt afgegeven, hangt niet af van wat er nog over is: een
gelijkmatige toevoer, en dat is precies het doel van een pleister.
\end{solution}

\begin{solution}{exo:b1:rate-laws:12}
$E_a = R\ln2/(1/298.15 - 1/308.15) = 8.314 \times 0.693/(1.088 \times
10^{-4}) = \qty{52.9}{kJ/mol}$. Tussen 273.15 en \qty{283.15}{K}:
$\exp\big(\frac{52900}{8.314}(\frac{1}{273.15} - \frac{1}{283.15})\big) =
2.3$: dezelfde \omterm{def:b1:rate-laws:activation-energy}{activeringsenergie} geeft bij lagere temperatuur een
grotere factor.
\end{solution}

\begin{solution}{pb:b1:rate-laws:1}
\textbf{1.} Verliezen van 8.6, 7.9, 7.2 en 6.5 procentpunt: niet
constant, dus niet van orde 0.
\textbf{2.} $\ln$: 4.605, 4.515, 4.425, 4.335, 4.245; het daalt met 0.090
per 10 dagen: een rechte, orde 1.
\textbf{3.} $k_{60} = 0.090/10 = \qty{9.0e-3}{d^{-1}}$.
\textbf{4.} $t_{1/2} = \ln2/k_{60} = \qty{77}{d}$.
\textbf{5.} $100\,\eu^{-0.54} = 58.3\,\%$.
\textbf{6.} $t_{90} = \ln(10/9)/k_{60} = 0.1054/0.0090 = \qty{11.7}{d}$.
\textbf{7.} Er blijft 90\,\% over: $\eu^{-kt_{90}} = 0.9$, dus
$t_{90} = \ln(10/9)/k$.
\textbf{8.} $t_{90}/t_{1/2} = \ln(10/9)/\ln2 = 0.152$.
\textbf{9.} Bij orde 1 hangt de resterende fractie alleen af van $kt$,
niet van de beginhoeveelheid.
\textbf{10.} Bij \qty{25}{\celsius} duurt de afbraak jaren; verwarmen
versnelt ze (Arrhenius), zodat ze in weken gemeten kan worden.
\textbf{11.} Bij orde 2 is de \omterm{def:b1:rate-laws:half-life}{halveringstijd} $1/(ak\,a_0)$: een tablet
die twee keer zo sterk is, zou relatief twee keer zo snel afbreken en
minder lang houdbaar zijn.
\textbf{12.} $1/T$ (\unit{K^{-1}}): \num{3.1934e-3}, \num{3.0945e-3},
\num{3.0017e-3}; $\ln k$: $-6.669$, $-5.661$, $-4.711$.
\textbf{13.} Richtingscoëfficiënt $= (-4.711 + 6.669)/(3.0017 \times 10^{-3} - 3.1934 \times
10^{-3}) = \qty{-1.021e4}{K}$.
\textbf{14.} $E_a = 8.314 \times 1.021 \times 10^4 = \qty{85}{kJ/mol}$.
\textbf{15.} Voorspeld: $\ln k_{50} = -4.711 - 1.021 \times 10^4 \times (3.0945 -
3.0017) \times 10^{-3} = -5.658$, gemeten $-5.661$: op de rechte.
\textbf{16.} $A = k_{60}\,\eu^{E_a/RT} = 0.0090\,\eu^{10215/333.15} =
\qty{1.9e11}{d^{-1}}$.
\textbf{17.} $\exp\big(1.021 \times 10^4(\frac{1}{298.15} -
\frac{1}{308.15})\big) = 3.0$: meer dan de factor 2 van de regel, omdat
$E_a$ groter is dan \qty{53}{kJ/mol}.
\textbf{18.} $k_{30} = k_{60}\exp\big(-1.021 \times 10^4(\frac{1}{303.15} -
\frac{1}{333.15})\big) = \qty{4.3e-4}{d^{-1}}$.
\textbf{19.} $t_{90} = 0.1054/4.33 \times 10^{-4} = \qty{243}{d}$, 8.0
maanden.
\textbf{20.} $1 - \eu^{-3 \times 0.0090} = 2.7\,\%$.
\textbf{21.} De houdbaarheid daalt snel met de temperatuur (8 maanden bij
\qty{30}{\celsius} tegen 14 bij \qty{25}{\celsius}): de vervaldatum geldt
alleen beneden \qty{25}{\celsius}.
\textbf{22.} $k_{25} = k_{60}\exp\big(-1.021 \times 10^4(\frac{1}{298.15} -
\frac{1}{333.15})\big) = \qty{2.46e-4}{d^{-1}}$.
\textbf{23.} $t_{90} = 0.1054/2.46 \times 10^{-4} = \qty{428}{d}$, dat is
\textbf{14 maanden}: de datum die op de doos gedrukt staat, verkregen uit
zes weken metingen.
\end{solution}
